\chapter{Formulação da cinemática inversa}
\label{ch:inv_kin}

A cinemática inversa permite obter os ângulos das juntas $(J1, \dots, J6)$ a partir da posição da garra $(x, y, z)$. 
Ao contrário da cinemática direta, a solução não é sempre única, pois o mesmo ponto pode ser obtido com o cotovelo para cima ou para baixo, e pode não existir, se o destino estiver fora do alcance do braço.

Seguindo a mesma abordagem da cinemática direta, o braço é decomposto em duas partes mais simples: um braço planar, que se move num plano vertical, e uma rotação desse plano em torno do eixo vertical.

\begin{itemize}
    \item No \textbf{plano horizontal}, a junta 1 roda o braço em torno do eixo vertical, logo define a orientação do plano vertical no qual o braço se movimenta;
    \item No \textbf{plano lateral}, as juntas 2, 3 e 4 comportam-se como um braço planar de 3 DOF, simplificando para 2 DOFs ao impor o ângulo da junta 4 para definir a orientação da garra. 
    \item A \textbf{junta 5} é fixada a $0^\circ$, e a \textbf{junta 6} roda a garra para alinhar com o objeto a ser agarrado.
\end{itemize}

% Secções do capítulo
\section{\textit{Plano horizontal}}
\label{sec:plano_horizontal}

Numa vista superior do braço, como é possível observar na Figura~\ref{fig:topo}, a \underline{junta 1} é a única junta que roda o plano do braço em torno do eixo vertical. 
As juntas 2, 3 e 4 não mudam a direção desse plano: mudam apenas a distância $d$ em frente e a altura $h$, calculadas na Secção~\ref{sec:plano_lateral}.

\begin{figure}[h]
  \centering
  \resizebox{0.6\textwidth}{!}{%Vista de cima do robô
\begin{tikzpicture}[scale=0.3,
                    % Estilos
                    link/.style={line width=2pt, line cap=round},
                    cota/.style={gray!70!black, thin}]

  % Coordenadas das juntas
    \coordinate (J1) at (0,0);
    \coordinate (J6) at (-4.7,21);     % centro de preensão (x, y)
    \coordinate (O5) at (5.6,3.2);

  % Base
    \draw[->] (-10,0) -- (12,0) node[below left] {$x$};
    \draw[->] (0,-3) -- (0,25) node[left] {$y$};
    \fill[gray!25] (0,0) circle (2.5);  \draw (0,0) circle (2.5);
    \node[below left] at (-2.2,-0.2) {J1};
  
  % Desvio lateral L5 e plano do braço
    \draw[link] (J1) -- node[pos=0.55, below right=-2pt] {$L_5$} (O5);
    \draw[link] (O5) -- node[pos=0.4, right=2pt] {$d$} (J6);
    \draw[cota, dashed] (O5) -- ($(O5)!1.1!(J6)$) node[above, black] {plano do braço};
    \draw[cota, dashed] (O5) -- (30:14);
    \fill (J6) circle (7pt); \node[left=4pt] at (J6) {$(x, y)$};
  
  % Distância
    \draw[cota, dashed] (J1) -- node[pos=0.8, left=1pt, black] {$r$} (J6);
  
  % Ângulos
    \draw[->, thick] (5,0) arc[start angle=0, end angle=30, radius=5] node[pos=0, below right=2pt and -3pt] {$q_1 > 0$};
    \draw[->, thin] (8.5,0) arc[start angle=0, end angle=102.6, radius=8.5] node[pos=0.27, right=2pt] {$\operatorname{atan2}(y, x)$};
    \draw[->, thin] (30:12) arc[start angle=30, end angle=102.6, radius=12] node[pos=0.3, above right=0pt] {$\operatorname{atan2}(d, L_5)$};
    \draw[thin] ($(O5)+(120:1.3)$) -- ($(O5)+(120:1.3)+(210:1.3)$) -- ($(O5)+(210:1.3)$);

  % Convenção de sinal (visto de cima)
    \draw[->, thick, red!70!black] (9,23) ++(-30:1.6) arc[start angle=-30, end angle=250, radius=1.6];
    \node[red!70!black] at (9,23) {$+$};
    \node[right=2pt, align=left, font=\small] at (10.8,23) {ângulo positivo:\\sentido anti-horário};

\end{tikzpicture}}
  \caption{Vista superior do braço no seu plano horizontal.}
  \label{fig:topo}
\end{figure}

\subsection{Desvio lateral}

No MyCobot280, o braço possui um desvio lateral $L_5$ de aproximadamente 6,5 cm, que causa que o plano do braço não intercate o eixo de rotação da junta 1.
O desvio lateral é representado na Figura~\ref{fig:topo} como a distância entre a projeção da junta 3 no plano horizontal e o eixo de rotação da junta 1. 
O desvio lateral é um fator importante a ser considerado no cálculo da cinemática direta, pois ele afeta a posição final da garra no espaço tridimensional.

\subsection{Rotação da base}

Com $q_1 = 0$ a frente do braço aponta para $+y$ e o desvio $L_5$ fica em $+x$. Ao rodar a base de $q_1$ roda-se estas duas direções em torno de $z$, e basta projetar cada parcela nos eixos $x$ e $y$:
 
\begin{itemize}
    \item o desvio lateral $l \ (=L_5)$ faz um ângulo $q_1$ com o eixo $x$, logo contribui com \underline{$l\cos q_1$} em $x$ e \underline{$l\sin q_1$} em $y$;
    \item a distância $d$ em frente é perpendicular ao desvio, ou seja, faz um ângulo $q_1 + 90^\circ$ com o eixo $x$, logo contribui com \underline{$d\cos(q_1 + 90^\circ) = -d\sin q_1$} em $x$ e \underline{$d\sin(q_1 + 90^\circ) = d\cos q_1$} em $y$.
\end{itemize}
 
Somando as duas parcelas, e juntando a altura $h$ do plano lateral:
\begin{equation}
    x = l\cos q_1 - d\sin q_1, \qquad
    y = l\sin q_1 + d\cos q_1, \qquad
    z = h
    \label{eq:xyz}
\end{equation}
 
Sem desvio lateral ($l = 0$), a Equação~\eqref{eq:xyz} reduz-se à forma, $x = d\cos(q_1 + 90^\circ)$ e $y = d\sin(q_1 + 90^\circ)$: a distância ao eixo multiplicada pelo cosseno e pelo seno do ângulo da base.

\subsection{Distância ao eixo da base}

Por o desvio $L_5$ ser perpendicular a $d$, a distância $r$ da garra ao eixo de J1, vista de cima, é
\begin{equation}
    r = \sqrt{x^2 + y^2} = \sqrt{l^2 + d^2}
    \label{eq:r}
\end{equation}
e a direção da garra é $\operatorname{atan2}(y, x) = q_1 + \operatorname{atan2}(d, l)$, como se vê na Figura~\ref{fig:topo}. Esta relação é a usada, ao contrário, na cinemática inversa para obter $q_1$.
 

\clearpage
\section{\textit{Plano lateral}}
\label{sec:plano_lateral}

No plano do braço, como é possível observar na Figura \ref{fig:lateral}, as coordenadas da garra são expressas nos eixos $d$ e $h$, e os ângulos são medidos a partir do eixo horizontal.   

\begin{figure}[h]
  \centering
  \resizebox{0.6\textwidth}{!}{% Cinemática inversa no plano do braço: ponto do pulso e braço de 2 DOF
% Esboço (pose ilustrativa); unidades em cm; origem no ombro (J2)
\begin{tikzpicture}[scale=0.32, >=stealth,
    elo/.style={line width=2.2pt, line cap=round},
    junta/.style={circle, draw, fill=white, line width=1pt, inner sep=2.2pt},
    cota/.style={gray!70!black, thin},
    ang/.style={->, thin}]
  
    % Coordenadas
    \coordinate (J2) at (0,0);
    \coordinate (J3) at (3.8,10.3);
    \coordinate (J4) at (13.2,8.7);
    \coordinate (J5) at (20.5,8.7);
    \coordinate (G)  at (20.5,-6.4);

  % Eixos
    \draw[->, thin] (-3,0) -- (25,0) node[below left] {$d$};
    \draw[->, thin] (0,-8) -- (0,15) node[left] {$h - L_1$};

  % Último elo
    \draw[line width=1.6pt, dashed, gray] (J4) -- node[below] {$L_4$} (J5) -- node[right] {$D$} (G);
    \fill (G) circle (6pt); \node[below right=1pt and 3pt] at (G) {alvo $(d,\, z - L_1)$};
    \draw[cota, dashed] (J5) -- ++(4.5,0);
    \draw[ang] ($(J5)+(3,0)$) arc[start angle=0, end angle=-90, radius=3] node[pos=0.5, right=2pt] {$\alpha_3$};

  % Braços
    \draw[elo] (J2) -- node[left] {$L_A = L_2$} (J3);
    \draw[elo] (J3) -- node[above] {$L_B = L_3$} (J4);
    \draw[cota, dashed] (J2) -- node[pos=0.6, below right=-1pt, black] {$r_3$} (J4);

  % Ângulos
    \draw[ang] (6,0) arc[start angle=0, end angle=33.4, radius=6] node[pos=0.5, right=1pt] {$\alpha_1$};
    \draw[ang] (33.4:4) arc[start angle=33.4, end angle=70, radius=4] node[pos=0.5, above right=0pt] {$\beta_1$};
    \draw[ang] ($(J3)+(249.8:2.5)$) arc[start angle=249.8, end angle=350.3, radius=2.5] node[pos=0.55, below=1pt] {$\beta_2$};

  % Juntas
    \node[junta] at (J2) {}; \node[below left=2pt] at (J2) {J2 (ombro)};
    \node[junta] at (J3) {}; \node[above=4pt] at (J3) {J3 (cotovelo)};
    \node[junta] at (J4) {}; \node[above right=3pt and -6pt] at (J4) {J4: $(x_{3d},\, y_{3d})$};
    \node[junta] at (J5) {}; \node[above=4pt] at (J5) {J5};
    
\end{tikzpicture}}
  \caption{Vista lateral do braço no seu plano vertical.}
  \label{fig:2dof}
\end{figure}

Obtendo as coordenadas da garra $(d, h)$ e os comprimentos dos braços $L_1$, $L_2$, $L_3$, $L_4$ e $D = L_6 + L_7$, é possível determinar os ângulos $\theta_A$, $\theta_B$, $\theta_C$ e $\alpha$.

Onde:

\begin{figure*}[!h]
\centering
    \begin{equation}
        \theta\textsubscript{A} = q\textsubscript{2} + \frac{\pi}{2}; \quad
        \theta\textsubscript{B} = \theta\textsubscript{A} + q\textsubscript{3}; \quad
        \theta\textsubscript{C} = \theta\textsubscript{B} + q\textsubscript{4}; \quad
        \alpha = \theta\textsubscript{C} + \frac{\pi}{2};
        \label{eq:2.1}
    \end{equation}
    \caption*{Equação (2.1)}
\end{figure*}

\clearpage



